\documentclass[a4paper]{article}
\usepackage{amsfonts}
\usepackage{amsmath}
\usepackage{amssymb}
\usepackage{graphicx}     

\begin{document}

\newcommand{\csch}{\operatorname{csch}}
\newcommand{\coth}{\operatorname{coth}}

\noindent \large Auteur: Sadia Nazary \\
\noindent \large Studierichting: Informatica \\
\noindent \large Oefeningenreeks 2F nr. 3.2 \\

\medskip

\normalsize

\noindent Bewijs dat: \\

$$
\coth^2 x - \csch^2 x = 1
$$

\noindent Uitgeschreven in termen van $ \sinh $ en $ \cosh $: \\

$$
\coth x = \dfrac{\cosh x}{\sinh x}, \quad \csch x = \dfrac{1}{\sinh x}
$$

$$
\Rightarrow \coth^2 x = \dfrac{\cosh^2 x}{\sinh^2 x}, \quad \csch^2 x = \dfrac{1}{\sinh^2 x}
$$

$$
\Rightarrow \coth^2 x - \csch^2 x = \dfrac{\cosh^2 x - 1}{\sinh^2 x}
$$

\noindent Gebruik de identiteit $ \cosh^2 x - \sinh^2 x = 1 \Rightarrow \cosh^2 x - 1 = \sinh^2 x $: \\

$$
= \dfrac{\sinh^2 x}{\sinh^2 x} = 1
$$

$$
\Rightarrow \coth^2 x - \csch^2 x = 1
$$

\begin{thebibliography}{999}

% \bibitem{a} Referentie

\end{thebibliography}

\end{document}
